\documentclass[12pt,a4paper]{article}

% Package yang diperlukan
\usepackage[bahasa]{babel}
\usepackage[utf8]{inputenc}
\usepackage{geometry}
\usepackage{graphicx}
\usepackage{setspace}
\usepackage{titlesec}
\usepackage{fancyhdr}

% Pengaturan margin: 4cm kiri, 3cm atas, 2.5cm kanan, 2.5cm bawah
\geometry{
    left=4cm,
    top=3cm,
    right=2.5cm,
    bottom=2.5cm
}

% Pengaturan spasi 1.5
\onehalfspacing

% Pengaturan font Times New Roman (gunakan mathptmx atau times)
\usepackage{mathptmx}

% Pengaturan header dan footer
\pagestyle{fancy}
\fancyhf{}
\fancyfoot[R]{\thepage}
\renewcommand{\headrulewidth}{0pt}

\begin{document}

% ==================== HALAMAN SAMPUL ====================
\begin{titlepage}
    \centering
    
    % Judul Proposal (gunakan huruf besar, tanpa singkatan)
    {\Large \textbf{JUDUL PROPOSAL TESIS ANDA}\par}
    
    \vspace{2cm}
    
    % Logo Universitas (sesuaikan path gambar)
    \includegraphics[width=5cm]{logo_unm.png}\par
    
    \vspace{1cm}
    
    {\Large \textbf{PROPOSAL TESIS}\par}
    
    \vspace{1cm}
    
    {\normalsize Diajukan sebagai salah satu syarat untuk memperoleh gelar\\
    Magister Komputer (M.Kom)\par}
    
    \vspace{2cm}
    
    {\normalsize
    \textbf{Nama Anda}\\
    NIM: 12345678\par}
    
    \vspace{2cm}
    
    {\normalsize
    \textbf{PROGRAM STUDI ILMU KOMPUTER (S2)}\\
    \textbf{FAKULTAS TEKNOLOGI INFORMASI}\\
    \textbf{UNIVERSITAS NUSA MANDIRI}\\
    \textbf{JAKARTA}\\
    \textbf{2024}\par}
    
\end{titlepage}

% ==================== HALAMAN PENGESAHAN ====================
\newpage
\thispagestyle{empty}

\begin{center}
    {\Large \textbf{HALAMAN PENGESAHAN}\par}
\end{center}

\vspace{1cm}

\noindent
Proposal Tesis ini diajukan oleh:\\

\noindent
\begin{tabular}{ll}
Nama & : Nama Anda \\
NIM & : 12345678 \\
Program Studi & : Ilmu Komputer \\
Judul & : Judul Proposal Tesis Anda \\
\end{tabular}

\vspace{2cm}

\noindent
Telah disetujui oleh:

\vspace{3cm}

\noindent
\begin{tabular}{p{7cm}p{7cm}}
Pembimbing, & Jakarta, \underline{\hspace{3cm}} 2024\\
& \\
& \\
& \\
\underline{\hspace{6cm}} & \\
Nama Dosen Pembimbing & \\
\end{tabular}

% ==================== BAB I PENDAHULUAN ====================
\newpage
\setcounter{page}{1}

\section{PENDAHULUAN}

\subsection{Latar Belakang}
Tuliskan latar belakang penelitian Anda di sini. Jelaskan:
\begin{itemize}
    \item Persoalan-persoalan yang muncul dan dihadapi
    \item Kesenjangan antara keadaan nyata dan keadaan ideal
    \item Inti masalah yang harus diselesaikan
    \item Rujukan pada penelitian sebelumnya dengan mencantumkan sumber
    \item Makna penting mengapa penelitian ini harus dilaksanakan
\end{itemize}

Contoh: Perkembangan teknologi informasi saat ini telah membawa perubahan signifikan dalam berbagai aspek kehidupan \cite{contoh2023}.

\subsection{Identifikasi Masalah}
Tuliskan identifikasi masalah secara ringkas dan terfokus. Ini merupakan intisari dari latar belakang.

Contoh format:
\begin{enumerate}
    \item Masalah pertama yang diidentifikasi
    \item Masalah kedua yang diidentifikasi
    \item Masalah ketiga yang diidentifikasi
\end{enumerate}

\subsection{Tujuan Penelitian}
Tuliskan tujuan penelitian Anda. Ini memuat hal-hal yang akan dicapai setelah penelitian berakhir.

Contoh:
Tujuan dari penelitian ini adalah:
\begin{enumerate}
    \item Tujuan pertama
    \item Tujuan kedua
    \item Tujuan ketiga
\end{enumerate}

\subsection{Ruang Lingkup Penelitian}
Tuliskan batasan penelitian Anda. Ini merupakan penegasan bagian masalah yang akan dipecahkan, termasuk asumsi-asumsi yang digunakan.

Ruang lingkup penelitian ini meliputi:
\begin{itemize}
    \item Batasan pertama
    \item Batasan kedua
    \item Asumsi yang digunakan
\end{itemize}

\subsection{Sistematika Penulisan}
Tuliskan sistematika penulisan proposal tesis Anda.

\textbf{BAB I PENDAHULUAN}\\
Bab ini berisi latar belakang, identifikasi masalah, tujuan penelitian, ruang lingkup penelitian, dan sistematika penulisan.

\textbf{BAB II LANDASAN/KERANGKA PEMIKIRAN}\\
Bab ini berisi kerangka teori dan tinjauan organisasi/obyek penelitian (opsional).

\textbf{BAB III METODOLOGI PENELITIAN}\\
Bab ini berisi metodologi yang akan digunakan dalam penelitian.

% ==================== BAB II LANDASAN PEMIKIRAN ====================
\newpage

\section{LANDASAN/KERANGKA PEMIKIRAN}

\subsection{Kerangka Teori}
Tuliskan teori-teori yang relevan, lengkap, mutakhir, dan urut sesuai dengan permasalahan penelitian.

\subsubsection{Teori Pertama}
Jelaskan teori pertama yang mendukung penelitian Anda \cite{referensi2024}.

\subsubsection{Teori Kedua}
Jelaskan teori kedua yang relevan dengan penelitian.

\subsection{Tinjauan Organisasi/Obyek Penelitian (Opsional)}
Jika penelitian Anda melibatkan organisasi atau obyek tertentu, tuliskan profil dan uraiannya di sini.

% ==================== BAB III METODOLOGI PENELITIAN ====================
\newpage

\section{METODOLOGI PENELITIAN}

Pada bab ini, jelaskan secara bertahap dan terperinci tentang:
\begin{itemize}
    \item Langkah-langkah (metodologi) yang akan digunakan
    \item Cara mencari dan mengumpulkan data
    \item Metode analisis yang akan digunakan
    \item Alat-alat, pendekatan, atau algoritma yang akan digunakan
    \item Uraian model yang akan digunakan
\end{itemize}

Metodologi ini harus merujuk kepada tujuan penelitian yang ditulis pada Bab I.

% Contoh menambahkan gambar metodologi
\begin{figure}[h]
    \centering
    \includegraphics[width=0.8\textwidth]{metodologi.png}
    \caption{Kerangka Metodologi Penelitian}
    \label{fig:metodologi}
\end{figure}

% ==================== DAFTAR REFERENSI ====================
\newpage

\begin{thebibliography}{99}

\bibitem{contoh2023}
Nama Penulis, ``Judul Artikel,'' \textit{Nama Jurnal}, vol. X, no. Y, pp. ZZ-ZZ, Tahun.

\bibitem{referensi2024}
Nama Penulis, \textit{Judul Buku}, Edisi, Penerbit, Kota, Tahun.

% Tambahkan referensi lainnya di sini
% Minimal 15 referensi ilmiah

\end{thebibliography}

\end{document}
